\section{General Tips for Success}
%=============================================================================

\subsection{Time Management}

\begin{itemize}[leftmargin=*, itemsep=4pt]
    \item \textbf{Front-load the hard parts.} Start with whatever you are most uncertain about. The biggest regret teams have at the end of the semester is ``we should have started building/coding/testing earlier.''
    \item \textbf{Respect lead times.} Components take time to ship. The machine shop has a queue. Resin prints can fail and need reprinting. Build buffer into every timeline.
    \item \textbf{Set internal deadlines 48 hours before the real ones.} This gives you time to integrate, review, and fix issues before submission. Last-minute scrambles produce sloppy work and missed details.
\end{itemize}

\subsection{Team Dynamics}

\begin{itemize}[leftmargin=*, itemsep=4pt]
    \item \textbf{Define roles and ownership early.} Every major subsystem and every deliverable should have one person who is ultimately responsible, even if others contribute.
    \item \textbf{Communicate proactively.} If you are stuck, behind schedule, or unsure about something, tell your team \textit{now}, not the night before a deadline. No one can help you with a problem they do not know about.
    \item \textbf{Use the retrospectives honestly.} They exist specifically to surface and address team process issues. If something is not working, the retrospective is the time and place to address it constructively.
    \item \textbf{Your individual grade depends on peer evaluations.} A peer-review-based adjustment factor is applied to the group portion of your grade. If your teammates consistently report that you are not contributing, your grade will be reduced---potentially significantly. Conversely, teammates who go above and beyond may receive a boost. This adjustment is \textit{not} visible in Canvas calculated grades; it is applied at the end of the semester.
    \item \textbf{Persistent non-contribution triggers a formal process.} Teammates who fail to contribute, communicate, or behave professionally may be placed on a Performance Improvement Plan (PIP) with one week to demonstrate improvement. If the issues persist, the student is removed from the team and must complete all remaining teamwork individually. Address problems early through retrospectives and direct conversation---do not let issues fester until someone is on a PIP.
\end{itemize}

\subsection{Calendar Awareness}

Several schedule disruptions are built into the \currentsemester{} semester. Account for them from the start:

\begin{itemize}[leftmargin=*, itemsep=4pt]
    \item \textbf{Career Day (\careerday, Tuesday):} No class. Falls during Sprint~2 (PDR) work. Plan your reverse engineering and FMEA work around this lost day.
    \item \textbf{Presidents' Day Break (\presidentsday, Tuesday):} No class. Falls during the Sprint~2/3 transition when you are finalizing your PDR and beginning CDR preparation.
    \item \textbf{Spring Break (\springbreakdates):} Campus is closed---no shipping, no machine shop, no lab access. All critical component orders and 3D print submissions should be complete with margin for iteration \textit{before} \springbreakorderby.
    \item \textbf{Attendance:} More than 3 unexcused absences results in automatic failure. Being physically present but not participating (playing games, watching videos, working on other classes) counts as absent.
\end{itemize}

\subsection{Documentation Habits}

\begin{itemize}[leftmargin=*, itemsep=4pt]
    \item \textbf{Document as you go.} Taking a photo of your prototype before and after every modification takes seconds and saves hours when writing the final report.
    \item \textbf{Version control your CAD and code.} Use Git (even GitHub Desktop if the command line is intimidating) for all firmware and software. For CAD files, use a consistent naming convention in a shared Google Drive folder: \texttt{Enclosure\_V3-2\_2026-04-15.SLDPRT} is infinitely better than \texttt{Final\_Final\_reallyFinal.SLDPRT}. Whatever system you choose, start it on Day~1, not the night before CDR.
    \item \textbf{Keep your BoM current.} Update it every time you order, swap, or discard a component. Reconstructing a BoM from memory at the end of the semester is painful and error-prone.
    \item \textbf{Document \textit{why}, not just \textit{what}.} The most perishable artifact in any engineering project is the rationale for design decisions. Record why you chose this sensor over that one, why you changed the mounting approach in Week~8, and what known limitations or workarounds exist. Your final report's ``Lessons Learned'' and ``Design Evolution'' sections will practically write themselves if you have been recording rationale all along.
    \item \textbf{Maintain a running Lessons Learned log.} As a team, spend 5 minutes at the end of every sprint adding to a shared document: ``What went well? What went wrong? What will we change?'' This single log will make writing your final ``Lessons Learned'' section a 15-minute task instead of a painful, hours-long reconstruction.
\end{itemize}

\subsection{Prototyping and Testing}

\begin{itemize}[leftmargin=*, itemsep=4pt]
    \item \textbf{Test early and test often.} Do not wait until Sprint~5 to see if your electronics and mechanical systems work together. Bench-test individual components as soon as they arrive.
    \item \textbf{Have a ``Plan B'' component.} For any single-source critical component, know what your fallback is if it is out of stock, damaged on arrival, or does not perform as specified.
    \item \textbf{Design for debugging.} Include test points in your circuits, print housings that can be opened, and write code with serial debug output. Your future self debugging at 10~PM will appreciate it.
    \item \textbf{Use keyed connectors everywhere.} If your design has multiple similar-looking electrical connections, specify keyed connectors that can only be inserted one way. Incorrect mating is one of the most common causes of electrical failure on first power-up in student projects, and it destroys components.
    \item \textbf{Never design your own AC-DC converter.} Always use a pre-certified (UL/CE listed) wall adapter. These are fully enclosed, safety-rated, and eliminate all mains-voltage hazards from your design. If your project involves lithium batteries, a Battery Management System (BMS) is mandatory for overcharge, deep-discharge, and short-circuit protection. Never charge lithium cells unattended; a LiPo fire releases toxic gas and can close a lab. Always charge in a fireproof LiPo safety bag. See MRD Chapter~12 (Power Architecture); search the index for ``lithium battery safety'' and ``BMS.''
    \item \textbf{If you design a custom PCB, print it at 1:1 scale and place real components on the printout before ordering.} This single step catches 80\% of footprint errors, the most expensive PCB mistake, because a wrong footprint means a scrapped board. Also: place a 100\,nF decoupling capacitor within 5\,mm of every IC power pin, and fill your bottom layer with a ground copper pour. These are not optional.
    \item \textbf{Use state machines and \texttt{millis()}-based timing in your firmware.} Avoid \texttt{delay()} calls, which block the CPU and make your system unresponsive. Structure your code around states (IDLE, READING, TRANSMITTING, SLEEPING) and use non-blocking timing. Enable the watchdog timer for crash recovery. Add a heartbeat LED as a system health indicator; if the LED stops blinking, something has hung.
\end{itemize}

\subsection{Using Campus Resources}

\begin{itemize}[leftmargin=*, itemsep=4pt]
    \item \textbf{Bulk wire, heat shrink, and perf boards} are available on campus for free or at minimal cost (see Appendix B of the project document). Wire is available in BBW160 from TAs during office hours: 18~AWG stranded at \$0.19/ft, 24~AWG stranded at \$0.08/ft. Solderable perf boards are stocked in CTLM123 in several sizes. Include these per-unit costs in your BoM for accurate accounting, but do not waste budget ordering them from vendors.
    \item \textbf{Get manufacturing reviews early.} The machine shop assistants, Prof.\ Blood, and Master TAs can catch design-for-manufacturing issues before you waste material and time. This review is required, but even if it were not, you should do it.
    \item \textbf{Office hours exist for a reason.} Your instructor and TAs have seen dozens of these projects. Five minutes of their input can save you days of going down the wrong path.
    \item \textbf{Use the Master Reference Document.} The \textit{Master Reference Document} is a comprehensive technical reference organized in three parts: Part~I covers the V-model design-build lifecycle with worked GreenShield examples, Part~II provides detailed electromechanical design guidance including connectors, power supplies, motors, power architecture, fasteners, bearings, power transmission, injection molding, PCB design, and microcontroller selection, and Part~III walks through PDR, CDR, and FDR deliverables with evaluation rubrics. Every Part~II chapter ends with a Requirements Mapping box showing how to translate technical specs into testable shall-statements. Consult the Quick Reference Tables appendix for consolidated design values.
\end{itemize}

%=============================================================================
